% ============================================================
% General Introduction
% ============================================================

\chapter*{General Introduction}
\addcontentsline{toc}{chapter}{General Introduction}
\label{chap:general-introduction}

% ---- Context ------------------------------------------------
\section*{Context}

Agricultural activity in Algeria's southern regions depends on a set of conditions
that differ fundamentally from those found in temperate farming environments. The
Saharan zone, which covers the vast majority of the country's land area, sustains
agricultural production through oasis-based cultivation and irrigated field plots
fed by underground water resources. In regions such as El~Oued, crops including
alfalfa are grown in conditions of extreme summer heat, scarce rainfall, and limited
access to persistent network infrastructure. Farmers who monitor several spatially
distributed plots must do so with limited real-time visibility into soil moisture
trends, atmospheric conditions, and the effectiveness of irrigation applications.

The growing availability of low-cost embedded hardware, long-range wireless
communication technologies, and cloud-based storage has opened a viable path toward
automated field monitoring in these environments. IoT-based agricultural monitoring
systems have been deployed across diverse contexts over the past decade, demonstrating
that continuous, structured data collection from remote sensor nodes is achievable
even under constrained infrastructure conditions. However, the deployment of such
systems in Saharan agricultural settings, particularly for alfalfa monitoring,
remains underexplored, and several recurring design limitations in existing systems
have not been addressed.

% ---- Problem Statement --------------------------------------
\section*{Problem Statement}

This work identifies five interconnected problems that motivate the proposed system.

First, agricultural IoT systems deployed in areas with unreliable Internet
connectivity cannot depend on continuous data streaming. Without a local storage
mechanism that operates independently of network availability, measurement data
acquired during connectivity outages is lost. This creates gaps in the historical
record that cannot be recovered retrospectively.

Second, the intermittent power supply common to Saharan field sites requires that
battery-powered sensor nodes operate for extended periods without replacement.
Existing systems rarely describe the engineering mechanisms by which they achieve
low-power operation, limiting the reproducibility of their designs in similarly
constrained environments.

Third, reviewed agricultural IoT systems frequently store or process data using
the time of upload to the server rather than the time of measurement as recorded
by the node. When connectivity is intermittent and upload sessions are batched, this
conflation introduces systematic temporal inaccuracy into all downstream analysis.
Charts, trend comparisons, and agronomic event alignments built on upload timestamps
are misleading when the gap between measurement and upload spans hours or days.

Fourth, no reviewed work targeting field IoT monitoring presents a deterministic
analytics layer that is separate from the visualization interface and explicitly
independent of any AI component. Without such a layer, quality control, reliability
scoring, and alert evaluation cannot be audited or reproduced from the data alone.

Fifth, despite its agronomic importance in Saharan agriculture, alfalfa has not
been the primary target crop in any of the reviewed IoT monitoring works. Systems
designed for generic crops or tropical contexts cannot be assumed to address the
monitoring requirements specific to alfalfa under the environmental conditions of
the El~Oued region.

% ---- Motivation ---------------------------------------------
\section*{Motivation}

This work is motivated by the intersection of two practical needs. The first is
the need for a reliable, traceable monitoring infrastructure for alfalfa cultivation
in the El~Oued region that operates within the constraints of intermittent power
and connectivity. The second is the academic need for a monitoring platform whose
data-handling decisions — particularly the separation of measurement time from
upload time — can be formally justified and validated through documented evidence.

The system developed in this thesis is designed to demonstrate that these two needs
can be addressed simultaneously through careful architectural choices in firmware,
data policy, backend design, and analytics pipeline structure, without requiring
commercial-grade infrastructure or continuous Internet access.

% ---- Objectives ---------------------------------------------
\section*{Objectives}

The work pursued the following objectives:

\begin{enumerate}
  \item Design a multi-node IoT architecture suited to Saharan field monitoring,
        with a central gateway coordinating remote sensing nodes over a LoRa
        wireless link.
  \item Implement structured soil and weather data collection across three physically
        distinct sensor nodes covering two irrigation pivots and one atmospheric
        monitoring position.
  \item Provide reliable local data storage on the gateway SD card that accumulates
        measurements independently of network availability.
  \item Design a manual, batch-triggered upload mechanism that transfers SD-stored
        records to a backend when Internet connectivity is temporarily available,
        without requiring continuous connectivity.
  \item Develop a backend and web dashboard that store, retrieve, and visualise sensor
        data using the RTC-derived measurement timestamp as the canonical time
        reference.
  \item Implement a deterministic quality-control and reliability-scoring analytics
        pipeline that operates entirely on processed data without AI involvement.
  \item Integrate a reliability-aware AI interpretation interface that provides
        narrative decision support from a controlled analytics snapshot, subject to
        explicit architectural constraints that prevent AI from computing primary
        metrics or modifying stored data.
  \item Validate the end-to-end system through a documented pilot deployment covering
        multi-node data acquisition, upload sessions, deterministic analytics
        output, and dashboard visualization.
\end{enumerate}

% ---- Contributions ------------------------------------------
\section*{Contributions}

The primary technical contributions of this work are:

\begin{enumerate}
  \item An offline-first, three-node LoRa IoT architecture with a structured
        ten-minute superframe that coordinates acquisition, reception, acknowledgement,
        and storage without any cloud dependency during field operation.
  \item A timestamp traceability policy centred on the DS3231 RTC, in which the
        \texttt{measured\_at} field is set once at packet reception and is never
        overwritten or approximated by the upload receipt timestamp.
  \item An SD-based local storage mechanism combined with a manual upload-on-demand
        trigger that makes the system fully operational without stable Internet
        access.
  \item Eight documented firmware operating modes — including Power Saving Mode,
        Recovery Mode for timing-drift correction, Upload Mode, and RTC
        Synchronisation Mode — each motivated by a specific observed field constraint.
  \item A ten-stage deterministic QC and reliability-scoring pipeline that produces
        auditable analytics snapshots from sensor readings without AI involvement at
        any processing stage.
  \item A Processed Data Viewer that exposes the full lineage of the analytics
        pipeline in read-only form, supporting transparency and academic traceability.
  \item A backend-proxied AI interpretation interface subject to six architectural
        boundary rules that confine Gemini to reading the deterministic snapshot and
        producing narrative text, without access to raw data or the ability to compute
        primary metrics.
  \item A documented prototype deployment targeting alfalfa monitoring in the
        El~Oued Saharan agricultural region, validated over a five-and-a-half-day
        pilot window.
\end{enumerate}

% ---- Thesis Organisation ------------------------------------
\section*{Thesis Organisation}

This thesis is organised into three chapters, preceded by this general introduction
and followed by a general conclusion.

\textbf{Chapter~1 — Background and Theoretical Foundations} establishes the context
and technical vocabulary needed to understand the design decisions in Chapter~3.
It introduces the Saharan agricultural environment that motivates the project,
discusses data collection principles for agricultural IoT systems, surveys wireless
communication technologies with a focus on LoRa, examines low-power and reliability
strategies for field-deployed embedded systems, and provides the agronomic context
of alfalfa cultivation. The chapter closes by identifying why alfalfa was selected
as the target crop and what monitoring requirements are specific to it.

\textbf{Chapter~2 — Related Works and Comparative Analysis} examines existing
agricultural IoT monitoring systems along seven dimensions: system architecture,
wireless communication, soil and environmental data collection, energy management
and reliability, offline storage and upload handling, timestamp semantics, and
dashboard plus analytics design. The chapter identifies alfalfa as absent from
the primary target crops of the reviewed works and synthesises five design gaps
that the proposed system is constructed to address.

\textbf{Chapter~3 — System Design, Implementation, and Evaluation} presents the
full implementation of the proposed platform, from hardware node architecture and
firmware operating modes through data storage policy, backend design, dashboard
implementation, deterministic analytics, AI interpretation boundary, and field
deployment. It concludes with a pilot validation section supported by quantitative
evidence from the deployment window and an explicit statement of system limitations.
